\documentclass[12pt,a4paper]{article}

\usepackage[utf8]{inputenc}
\usepackage[T1]{fontenc}
\usepackage[spanish,es-tabla]{babel}
\usepackage{geometry}
\usepackage{graphicx}
\usepackage{float}
\usepackage{booktabs}
\usepackage{longtable}
\usepackage{array}
\usepackage{enumitem}
\usepackage{xcolor}
\usepackage{hyperref}
\usepackage{titlesec}
\usepackage{tabularx}
\usepackage[table]{xcolor}
\usepackage{parskip}
\usepackage{caption}
\usepackage{listings}
\usepackage{multirow}

\geometry{margin=2.5cm}
\setlength{\parskip}{0.6em}
\setlength{\parindent}{0pt}

\hypersetup{colorlinks=true, linkcolor=blue, urlcolor=blue, citecolor=blue}

\titleformat{\section}{\normalfont\Large\bfseries}{\thesection.}{0.5em}{}
\titleformat{\subsection}{\normalfont\large\bfseries}{\thesubsection.}{0.5em}{}
\titleformat{\subsubsection}{\normalfont\normalsize\bfseries}{\thesubsubsection.}{0.5em}{}

\definecolor{codebg}{rgb}{0.96,0.96,0.96}
\definecolor{codekw}{rgb}{0.0,0.0,0.6}
\definecolor{codecomment}{rgb}{0.0,0.5,0.0}
\lstdefinestyle{sql}{
  language=SQL,
  backgroundcolor=\color{codebg},
  basicstyle=\ttfamily\footnotesize,
  keywordstyle=\color{codekw}\bfseries,
  commentstyle=\color{codecomment}\itshape,
  breaklines=true,
  frame=single,
  rulecolor=\color{gray!40},
  showstringspaces=false,
  morekeywords={EXPLAIN,ANALYZE,BUFFERS,BIGINT,VARCHAR,INTEGER,BETWEEN,COALESCE,NUMERIC,DEFERRABLE,DEFERRED},
  literate={á}{{\'{a}}}1 {é}{{\'{e}}}1 {í}{{\'{i}}}1 {ó}{{\'{o}}}1 {ú}{{\'{u}}}1
           {ñ}{{\~{n}}}1 {Á}{{\'{A}}}1 {É}{{\'{E}}}1 {Ó}{{\'{O}}}1 {Ú}{{\'{U}}}1 {Ñ}{{\~{N}}}1
}
\lstset{style=sql}

\begin{document}

% ------------------------------------------------------------------
% PORTADA
% ------------------------------------------------------------------
\thispagestyle{empty}
\begin{center}
\vspace*{1cm}
{\Large \bfseries Universidad de Ingeniería y Tecnología --- UTEC}\\[1.2cm]
{\Large \bfseries Proyecto de Base de Datos I}\\[0.3cm]
{\large Hito 2: Implementación en PostgreSQL, Optimización y Experimentación con Índices}\\[1.2cm]

\textbf{Curso:} Base de Datos I\\
\textbf{Profesor:} Wilder Nina Choquehuayta\\   % <-- confirmar con el profe real
\textbf{Tema:} Plataforma de catálogo, lectura, reseñas y descubrimiento de materiales literarios\\[1.2cm]

\textbf{Integrantes:}\\[0.2cm]
\textit{Sergio Felipe Delgado Amado --- 202310227}\\
\textit{Antonio Jorge Guerrero Cedrón --- 202510717}\\
\textit{Luciana Mylene Melgarejo Quispe --- 202420055}\\[1.2cm]

\textbf{Repositorio:}
\href{https://github.com/SefedeamU/Proyecto-BD-Hito-2}{github.com/SefedeamU/Proyecto-BD-Hito-2}

\vfill
{\large Ciclo 2026.1}
\end{center}

\newpage
\tableofcontents
\newpage

% ------------------------------------------------------------------
\section*{Sobre este informe}
\addcontentsline{toc}{section}{Sobre este informe}

Este documento corresponde al Hito 2 del proyecto. El Hito 1 (requisitos,
modelo E-R y modelo relacional) se entregó por separado y no se reproduce
aquí. El esquema implementado parte directamente de ese modelo; los ajustes
que fue necesario hacer al pasar al código se explican en la
Sección~\ref{sec:cambios} y están todos documentados en
\texttt{docs/cambios\_respecto\_hito1.md} dentro del repositorio.

El dominio es una plataforma de catálogo literario. La entidad central es
\texttt{Material}, que representa libros, ensayos, revistas, poemas y
audiolibros, y se relaciona con autores, géneros, subgéneros, editoriales,
premios, ilustraciones, curiosidades, reseñas, likes e historial de lectura.

% ==================================================================
\section{Implementación de la base de datos}

\subsection{Creación de tablas en PostgreSQL}

El esquema se implementó en PostgreSQL a partir del modelo relacional del
Hito 1. El script \texttt{sql/01\_tables.sql} crea las 25 tablas en el orden
correcto de dependencias. Para la traducción de tipos se siguió el diccionario
de datos del Hito 1:

\begin{center}
\begin{tabular}{ll}
\toprule
\textbf{Tipo en el modelo relacional} & \textbf{Tipo en PostgreSQL} \\
\midrule
Bigint           & \texttt{BIGINT} \\
Int              & \texttt{INTEGER} \\
Varchar(n)       & \texttt{VARCHAR(n)} \\
Date             & \texttt{DATE} \\
Double Precision & \texttt{NUMERIC(3,1)} (ver Sección~\ref{sec:cambios}, cambio G) \\
\bottomrule
\end{tabular}
\end{center}

Las restricciones de integridad del Hito 1 se implementaron de tres formas
según su naturaleza:

\begin{itemize}
  \item \textbf{Restricciones de columna o tabla} (\texttt{NOT NULL},
  \texttt{UNIQUE}, \texttt{PRIMARY KEY}, \texttt{FOREIGN KEY}, \texttt{CHECK}):
  para reglas que dependen solo de una fila o de una tabla. Por ejemplo,
  \texttt{CHECK (CHAR\_LENGTH(Password) = 12)} y \texttt{CHECK (Edad >= 12)}
  en \texttt{Usuario}; \texttt{CHECK (Relevancia BETWEEN 1 AND 5)} en
  \texttt{Premio}; \texttt{CHECK (... BETWEEN 1 AND 10)} para los cuatro
  indicadores de contenido en \texttt{AgeRate}.

  \item \textbf{Columna discriminadora}: la herencia total y disjunta de
  \texttt{Material} se garantiza con la columna \texttt{Tipo VARCHAR(10) NOT
  NULL} y un \texttt{CHECK} que solo acepta los cinco valores válidos
  (ver cambio A).

  \item \textbf{Triggers}: para reglas que dependen de dos tablas a la vez
  y no se pueden expresar con un \texttt{CHECK} simple
  (ver Sección~\ref{sec:triggers}).
\end{itemize}

Las definiciones de \texttt{Material} y \texttt{Usuario}, que concentran la
mayoría de las restricciones, son las siguientes:

\begin{lstlisting}
CREATE TABLE Material (
    Id               BIGINT       NOT NULL,
    Numero_Paginas   INTEGER      NOT NULL,
    Anio_Publicacion BIGINT       NOT NULL,
    Eslogan          VARCHAR(100),
    Alias            VARCHAR(50),
    Pais             VARCHAR(15)  NOT NULL,
    Idioma           VARCHAR(15)  NOT NULL,
    Tipo             VARCHAR(10)  NOT NULL,   -- discriminador de subtipo [cambio A]
    AgeRate_Code     BIGINT       NOT NULL,
    Editorial_Id     BIGINT       NOT NULL,
    CONSTRAINT pk_material      PRIMARY KEY (Id),
    CONSTRAINT ck_material_tipo CHECK (Tipo IN ('Libro','Ensayo','Revista','Poema','AudioBook')),
    CONSTRAINT fk_material_agerate   FOREIGN KEY (AgeRate_Code)  REFERENCES AgeRate (Code),
    CONSTRAINT fk_material_editorial FOREIGN KEY (Editorial_Id)  REFERENCES Editorial (Id)
);

CREATE TABLE Usuario (
    Username  VARCHAR(50)   NOT NULL,
    Email     VARCHAR(100)  NOT NULL,
    Password  VARCHAR(12)   NOT NULL,
    Nombre    VARCHAR(20)   NOT NULL,
    Apellido  VARCHAR(20)   NOT NULL,
    Edad      INTEGER       NOT NULL,
    Telefono  BIGINT        NOT NULL,
    Ciudad    VARCHAR(20)   NOT NULL,
    Rol       VARCHAR(20)   NOT NULL,
    CONSTRAINT pk_usuario          PRIMARY KEY (Username, Email),
    CONSTRAINT uq_usuario_username UNIQUE (Username),
    CONSTRAINT uq_usuario_email    UNIQUE (Email),
    CONSTRAINT ck_usuario_password CHECK (CHAR_LENGTH(Password) = 12),
    CONSTRAINT ck_usuario_edad     CHECK (Edad >= 12)
);
\end{lstlisting}

El esquema completo incluye las tablas maestras (\texttt{AgeRate},
\texttt{Editorial}, \texttt{Genero}, \texttt{Autor}, \texttt{Premio},
\texttt{Ilustracion}, \texttt{Curiosidad}, \texttt{Material},
\texttt{Usuario}), las entidades débiles (\texttt{Warning},
\texttt{SubGenero}), los cinco subtipos de Material y las tablas de relación
(\texttt{Escribe}, \texttt{Pertenece}, \texttt{PerteneceSubGenero},
\texttt{Popularizo}, \texttt{Contiene}, \texttt{Tiene}, \texttt{Ganar},
\texttt{Leer}, \texttt{Likes}, \texttt{Resena}).

\subsubsection{Triggers}
\label{sec:triggers}

El script \texttt{sql/02\_triggers.sql} implementa cuatro reglas del Hito 1
que dependen de más de una tabla:

\begin{enumerate}[label=(\alph*)]
  \item \textbf{Año de publicación vs.\ fundación editorial.}
  \texttt{Anio\_Publicacion} de un material debe ser mayor o igual al año de
  \texttt{Fundacion} de su editorial. Se valida con un trigger
  \texttt{BEFORE INSERT OR UPDATE} sobre \texttt{Material}.

  \item \textbf{Fecha de reseña vs.\ publicación.}
  El año de la fecha de una reseña debe ser mayor o igual al
  \texttt{Anio\_Publicacion} del material reseñado. También
  \texttt{BEFORE INSERT OR UPDATE}, esta vez sobre \texttt{Resena}.

  \item \textbf{Todo material tiene al menos un autor.}
  Participación total del lado \texttt{Material} en \texttt{Escribe}.
  Se implementa como \texttt{CONSTRAINT TRIGGER ... DEFERRABLE INITIALLY
  DEFERRED}: la validación ocurre al \texttt{COMMIT}, no fila por fila.
  Esto es necesario porque la inserción de \texttt{Material} y la de
  \texttt{Escribe} no se pueden hacer en el mismo instante; el trigger
  diferido da margen hasta que se cierre la transacción.

  \item \textbf{Todo género tiene al menos un autor representante.}
  Participación total en \texttt{Popularizo}. Mismo mecanismo diferido
  que el punto anterior.
\end{enumerate}

\subsubsection{Vistas de usuario}

El script \texttt{sql/03\_views.sql} define cinco vistas:

\begin{itemize}
  \item \textbf{vista\_material\_basico}: datos básicos de cada material
  con su editorial y clasificación de edad.
  \item \textbf{vista\_material\_autores}: materiales con sus autores
  (relación \texttt{Escribe}).
  \item \textbf{vista\_material\_subgenero}: materiales con género y
  subgénero; requiere la tabla \texttt{PerteneceSubGenero} añadida en
  el Hito 2 (cambio I).
  \item \textbf{vista\_popularidad\_material}: cantidad de likes, lecturas
  y reseñas por material, junto con el promedio de puntaje.
  \item \textbf{vista\_historial\_usuario}: historial de lectura por usuario
  con las fechas.
\end{itemize}

\subsection{Carga de datos}
\label{sec:carga}

La carga se realiza con un generador de datos sintéticos propio, en la carpeta
\texttt{faker/}. Está escrito en Python usando las librerías \texttt{Faker} y
\texttt{psycopg}. El generador aplica el esquema completo antes de poblar,
genera primero las tablas maestras y luego las de interacción para respetar las
llaves foráneas, y envuelve la carga de \texttt{Material}+\texttt{Escribe} y
\texttt{Genero}+\texttt{Popularizo} en transacciones para que los triggers
diferidos puedan validar al \texttt{COMMIT}. Usa \texttt{COPY} en lugar de
\texttt{INSERT} fila por fila para que la carga sea eficiente, y una semilla
fija (\texttt{SEED=42}) para que los datos sean reproducibles.

\subsection{Los cuatro escenarios de volumen}

Para el experimento se trabajó con cuatro bases de datos, una por escenario:

\begin{center}
\begin{tabular}{lrrrrr}
\toprule
\textbf{Base} & \textbf{Usuarios} & \textbf{Materiales}
  & \textbf{Likes} & \textbf{Lecturas} & \textbf{Reseñas} \\
\midrule
\texttt{bd\_literaria\_1k}   &     200 &     300
  &     1\,000 &     1\,000 &     1\,000 \\
\texttt{bd\_literaria\_10k}  &   1\,000 &   2\,000
  &    10\,000 &    10\,000 &    10\,000 \\
\texttt{bd\_literaria\_100k} &  10\,000 &  20\,000
  &   100\,000 &   100\,000 &   100\,000 \\
\texttt{bd\_literaria\_1m}   & 100\,000 & 100\,000
  & 1\,000\,000 & 1\,000\,000 & 1\,000\,000 \\
\bottomrule
\end{tabular}
\end{center}

La base de 1M tiene en total aproximadamente 3.86 millones de filas sumando
todas las tablas. Los cuatro escenarios se entregan como dumps restaurables en
\texttt{dumps/} y pueden regenerarse con el faker usando la misma semilla.

\subsection{Ajustes al modelo del Hito 1}
\label{sec:cambios}

Al escribir el código se identificaron trece puntos donde el modelo relacional
dejaba reglas del dominio sin mecanismo de cumplimiento, o donde había
decisiones de diseño que convenía precisar antes de poblar las bases. Todos
los cambios están marcados como comentarios \texttt{[Cambio X]} en los scripts
y explicados con detalle en \texttt{docs/cambios\_respecto\_hito1.md}.

\begin{center}
\begin{longtable}{clp{8.5cm}}
\toprule
\textbf{ID} & \textbf{Tipo} & \textbf{Descripción} \\
\midrule
\endhead
A & Integridad & Columna \texttt{Tipo} en \texttt{Material} para garantizar
    herencia total y disjunta. \\
B & Regla de negocio & Trigger diferido: un material sin autores no puede
    quedar en la base. \\
C & Regla de negocio & Trigger diferido: un género sin al menos un autor
    representante no puede quedar en la base. \\
D & Modelo & \texttt{Fecha} entra a la PK de \texttt{Leer} para poder
    registrar relecturas. \\
E & Integridad & \texttt{UNIQUE(Material\_Id, Usuario\_Username,
    Usuario\_Email)} en \texttt{Resena}: un usuario no puede reseñar el
    mismo material dos veces. \\
F & Integridad & \texttt{CHECK (Likes >= 0)} en \texttt{Resena}: el contador
    de likes no puede ser negativo. \\
G & Tipo de dato & \texttt{Resena.Puntaje} cambia de \texttt{DOUBLE PRECISION}
    a \texttt{NUMERIC(3,1)}: el Hito 1 especificaba un decimal exacto, y
    \texttt{DOUBLE PRECISION} no lo garantiza. \\
H & Integridad ref. & \texttt{ON UPDATE CASCADE} en las FK que apuntan a
    \texttt{Usuario}: si cambia el username o el email, las tablas dependientes
    se actualizan solas. \\
I & Modelo & Nueva tabla \texttt{PerteneceSubGenero}: el Hito 1 pedía poder
    buscar materiales por subgénero, pero el modelo relacional dejaba a
    \texttt{SubGenero} sin conexión con \texttt{Material}. \\
J & Experimento & La consulta de materiales populares se reescribió para evitar
    un producto cartesiano que distorsionaba los tiempos. \\
K & Experimento & Se quitaron los índices redundantes con la primera columna de
    alguna PK (no aportan mejora y ensucian el experimento). \\
L & Experimento & Nueva Consulta~5 que filtra por \emph{rango} de años
    (\texttt{BETWEEN}), para demostrar un índice B-tree sobre un predicado de
    rango (a pedido del profesor). \\
M & Modelo & Nueva tabla \texttt{ImagenMaterial}: cada material recibe un
    arreglo de tres URLs de portada (Lorem Picsum) para la aplicación de
    demostración. Los datos de imagen viven en la base, no en el front-end. \\
\bottomrule
\end{longtable}
\end{center}

Vale la pena detenerse en los cambios J y K porque afectan directamente los
resultados del experimento. El cambio J corrige la consulta de popularidad:
la versión original unía las tres tablas de interacción con \texttt{LEFT JOIN}
directos sobre \texttt{Material}, lo que generaba un producto cartesiano
(filas de likes por filas de lecturas por filas de reseñas por cada material).
La versión corregida agrega cada tabla por separado en subconsultas y luego las
une, de modo que el tiempo que se mide refleja el costo real de una agregación
total y no un error de formulación. El cambio K retira índices como el de
\texttt{Material\_Id} en \texttt{Escribe}: como esa columna es la primera de
la PK de \texttt{Escribe}, ya existe un índice de PK sobre ella y crear uno
secundario encima no cambia nada; incluirlo en el experimento solo añadiría
una fila de resultados con 0\% de mejora.

% ==================================================================
\section{Optimización y Experimentación}

\subsection{Consultas del experimento}

\subsubsection{Criterios de selección}

Se eligieron cinco consultas que cubren patrones de acceso distintos. El
criterio principal fue que el conjunto en su totalidad mostrara los tres
comportamientos posibles de un índice secundario: cuándo ayuda mucho,
cuándo ayuda poco y cuándo no ayuda. Una de las consultas es por rango
(\texttt{BETWEEN}), que es el caso más ilustrativo del funcionamiento de
un índice B-tree.

\begin{center}
\begin{longtable}{p{0.9cm}p{4.2cm}p{4.0cm}p{3.4cm}}
\toprule
\textbf{ID} & \textbf{Qué devuelve} & \textbf{Tipo de acceso}
  & \textbf{Índice relevante} \\
\midrule
\endhead
Q1 & Materiales de un autor & Filtro selectivo por igualdad
  & \texttt{idx\_escribe\_autor} \\
Q2 & Materiales de un género & Filtro poco selectivo por igualdad
  & \texttt{idx\_pertenece\_genero} \\
Q3 & Los 20 materiales más populares & Agregación total sin filtro
  & (ninguno aplica) \\
Q4 & Historial de lectura de un usuario & Igualdad + orden por fecha
  & \texttt{idx\_leer\_usuario} \\
Q5 & Materiales en un rango de años & \textbf{Rango} (\texttt{BETWEEN})
  & \texttt{idx\_material\_anio} \\
\bottomrule
\end{longtable}
\end{center}

\subsubsection{Consultas SQL}

Las consultas están limpias en \texttt{querys/} y con \texttt{EXPLAIN
(ANALYZE, BUFFERS)} en \texttt{sql/06\_experiments.sql}. Se reproducen
aquí las dos más representativas del experimento.

\paragraph{Q5 --- Materiales por rango de años.}

\begin{lstlisting}
SELECT m.Id, m.Alias, m.Anio_Publicacion, m.Idioma
FROM   Material m
WHERE  m.Anio_Publicacion BETWEEN 1900 AND 1905
ORDER  BY m.Anio_Publicacion;
\end{lstlisting}

El rango elegido (primeros años del catálogo) es selectivo: recupera
aproximadamente el 0.1\% de las filas en la base de 1M. Eso permite que
el planificador use el índice \texttt{idx\_material\_anio} de forma
natural, sin forzarlo. En la Sección~\ref{sec:crossover} se analiza
qué pasa cuando el rango es más amplio.

\paragraph{Q4 --- Historial de lectura de un usuario.}

\begin{lstlisting}
SELECT m.Id, m.Alias, le.Fecha
FROM   Usuario  u
JOIN   Leer     le  ON le.Usuario_Username = u.Username
                   AND le.Usuario_Email    = u.Email
JOIN   Material m   ON m.Id = le.Material_Id
WHERE  u.Username = :username
  AND  u.Email    = :email
ORDER  BY le.Fecha DESC;
\end{lstlisting}

El índice \texttt{idx\_leer\_usuario} es compuesto:
\texttt{(Usuario\_Username, Usuario\_Email, Fecha DESC)}. Además de
localizar las filas del usuario, entrega los resultados ya ordenados
por fecha descendente, lo que evita el paso de \texttt{Sort}.

\subsection{Metodología}

\begin{itemize}
  \item \textbf{Herramienta de medición:} \texttt{EXPLAIN (ANALYZE,
  BUFFERS)}, que ejecuta la consulta y reporta el tiempo real y el plan
  elegido por el planificador.

  \item \textbf{Número de corridas:} seis por consulta. La primera se
  descarta (estado frío) y se trabaja con la mediana de las cinco restantes
  para tener una medida estable en estado caliente.

  \item \textbf{Limpieza de caché entre corridas:} \texttt{DISCARD ALL}
  antes de cada corrida, para que el planificador vuelva a construir el
  plan desde cero y no reutilice planes cacheados de corridas anteriores.

  \item \textbf{Escenario sin índices:} se ejecuta \texttt{sql/05\_drop\_indexes.sql},
  que elimina todos los índices secundarios. Los de \texttt{PRIMARY KEY} y
  \texttt{UNIQUE} no se pueden quitar porque son parte del esquema; por eso
  las consultas del experimento filtran por columnas que no son la primera
  de ninguna PK (\texttt{Autor\_Id}, \texttt{Genero\_Nombre},
  \texttt{Usuario\_Username}/\texttt{Email}, \texttt{Anio\_Publicacion})
  o son agregaciones totales donde ningún índice aplica (Q3).

  \item \textbf{Escenario con índices:} se ejecuta \texttt{sql/04\_indexes.sql}
  y luego \texttt{ANALYZE} para que el planificador tenga estadísticas
  actualizadas.

  \item El protocolo se repite para las cuatro bases.
\end{itemize}

\subsection{Índices del experimento}

\begin{center}
\begin{tabular}{ll}
\toprule
\textbf{Nombre} & \textbf{Definición} \\
\midrule
\texttt{idx\_escribe\_autor}    & \texttt{Escribe(Autor\_Id)} \\
\texttt{idx\_pertenece\_genero} & \texttt{Pertenece(Genero\_Nombre)} \\
\texttt{idx\_resena\_material}  & \texttt{Resena(Material\_Id)} \\
\texttt{idx\_leer\_usuario}     & \texttt{Leer(Usuario\_Username, Usuario\_Email, Fecha DESC)} \\
\texttt{idx\_material\_anio}    & \texttt{Material(Anio\_Publicacion)} \\
\bottomrule
\end{tabular}
\end{center}

\subsection{Plataforma de pruebas}

Las mediciones se realizaron sobre PostgreSQL~16 en Ubuntu~24.04,
con la configuración por defecto del servidor, en las cuatro bases
descritas en la Sección~1.3.

\subsection{Medición de tiempos}

% ===========================================================================
% PENDIENTE — Completar tras correr el experimento en el servidor desplegado.
%
% Instrucciones para Sergio (o quien tenga acceso al servidor):
%
%   Para cada base (bd_literaria_1k, 10k, 100k, 1m):
%
%   1. Conectarse:
%      export PGPASSWORD='<TU_PASSWORD>'   # tomar de faker/.env (no versionado)
%      psql -h <HOST> -p 5432 -U <USUARIO> -d bd_literaria_1m
%
%   2. Quitar índices secundarios:
%      \i sql/05_drop_indexes.sql
%      ANALYZE;
%
%   3. Correr las consultas y anotar "Execution Time" de cada una:
%      \i sql/06_experiments.sql
%
%   4. Crear índices:
%      \i sql/04_indexes.sql
%      ANALYZE;
%
%   5. Correr de nuevo y anotar los nuevos "Execution Time":
%      \i sql/06_experiments.sql
%
%   Con esos 40 valores (5 consultas × 4 bases × 2 escenarios) se completan
%   las dos tablas de abajo, se generan los gráficos y se redactan los
%   análisis que aparecen como pendientes en este documento.
% ===========================================================================

\subsubsection{Sin índices}

\begin{center}
\begin{tabular}{lrrrr}
\toprule
\textbf{Consulta} & \textbf{1K (ms)} & \textbf{10K (ms)}
  & \textbf{100K (ms)} & \textbf{1M (ms)} \\
\midrule
Q1 (autor)     & \textit{pendiente} & \textit{pendiente} & \textit{pendiente} & \textit{pendiente} \\
Q2 (género)    & \textit{pendiente} & \textit{pendiente} & \textit{pendiente} & \textit{pendiente} \\
Q3 (populares) & \textit{pendiente} & \textit{pendiente} & \textit{pendiente} & \textit{pendiente} \\
Q4 (historial) & \textit{pendiente} & \textit{pendiente} & \textit{pendiente} & \textit{pendiente} \\
Q5 (rango)     & \textit{pendiente} & \textit{pendiente} & \textit{pendiente} & \textit{pendiente} \\
\bottomrule
\end{tabular}
\end{center}

\subsubsection{Con índices}

\begin{center}
\begin{tabular}{lrrrr}
\toprule
\textbf{Consulta} & \textbf{1K (ms)} & \textbf{10K (ms)}
  & \textbf{100K (ms)} & \textbf{1M (ms)} \\
\midrule
Q1 (autor)     & \textit{pendiente} & \textit{pendiente} & \textit{pendiente} & \textit{pendiente} \\
Q2 (género)    & \textit{pendiente} & \textit{pendiente} & \textit{pendiente} & \textit{pendiente} \\
Q3 (populares) & \textit{pendiente} & \textit{pendiente} & \textit{pendiente} & \textit{pendiente} \\
Q4 (historial) & \textit{pendiente} & \textit{pendiente} & \textit{pendiente} & \textit{pendiente} \\
Q5 (rango)     & \textit{pendiente} & \textit{pendiente} & \textit{pendiente} & \textit{pendiente} \\
\bottomrule
\end{tabular}
\end{center}

\subsection{Resultados}

% PENDIENTE — Una vez completadas las tablas de tiempos, esta sección debe
% incluir:
%   - Un gráfico de barras por consulta (sin vs. con índice, en las 4 bases).
%   - Un gráfico de aceleración (speedup) en la base de 1M.
%   - Análisis por consulta explicando qué plan usa el motor y por qué.
%
% Para cada consulta conviene capturar también la salida de:
%   EXPLAIN (ANALYZE, BUFFERS) <consulta>
% tanto sin índice como con índice, para incluirla como evidencia.

\textit{[Sección pendiente de completar una vez obtenidos los tiempos del
servidor.]}

\subsubsection{Q1 --- Materiales de un autor}

\textit{[Pendiente. Incluir: tiempos, plan EXPLAIN sin y con índice,
explicación de por qué \texttt{idx\_escribe\_autor} ayuda.]}

\subsubsection{Q2 --- Materiales de un género}

\textit{[Pendiente. Incluir: tiempos, plan EXPLAIN sin y con índice,
explicación del impacto de la selectividad del filtro.]}

\subsubsection{Q3 --- Materiales más populares}

\textit{[Pendiente. Incluir: tiempos, plan EXPLAIN sin y con índice,
explicación de por qué el índice no aporta en agregaciones totales.]}

\subsubsection{Q4 --- Historial de lectura de un usuario}

\textit{[Pendiente. Incluir: tiempos, plan EXPLAIN sin y con índice,
explicación de \texttt{idx\_leer\_usuario} y cómo evita el paso de Sort.]}

\subsubsection{Q5 --- Consulta por rango}

\textit{[Pendiente. Incluir: tiempos, plan EXPLAIN sin y con índice,
explicación del Bitmap Index Scan y la bondad del índice B-tree en
predicados BETWEEN selectivos. También analizar qué ocurre con rangos
amplios (el punto de cruce por selectividad).]}

\subsection{El costo de los índices}

% PENDIENTE — Correr sql/07_index_overhead.sql en la base de 1M y anotar:
%   - Tiempo de INSERT masivo sin índices.
%   - Tiempo de INSERT masivo con índices.
%   - Comparación de lectura no selectiva con y sin índice (forzado).

\textit{[Pendiente. El script \texttt{sql/07\_index\_overhead.sql} mide el
costo en escrituras y en lecturas no selectivas. Ejecutarlo sobre
\texttt{bd\_literaria\_1m} y completar con los resultados.]}

\subsection{Escalamiento por encima del millón (punto extra)}

\textit{[Pendiente. Una vez obtenidos los tiempos de las cuatro bases, graficar
el crecimiento del costo sin índices para ver si el patrón es lineal, y comparar
con el comportamiento con índices. Discutir qué estrategias serían necesarias
para volúmenes mayores al millón (vistas materializadas, particionamiento,
réplicas de lectura).]}

\subsection{Análisis y discusión}

\textit{[Pendiente. Una vez con los datos completos, redactar el análisis
comparativo de los tres comportamientos observados: cuándo el índice ayuda mucho,
cuándo poco y cuándo no ayuda. Incluir la aceleración (speedup) de cada consulta
en la base de 1M como referencia principal.]}

% ==================================================================
\section{Conclusiones}

\begin{itemize}
  \item Se implementó en PostgreSQL el modelo relacional del Hito 1: 25 tablas
  con restricciones de integridad (PK, FK, UNIQUE, NOT NULL, CHECK), cuatro
  triggers para reglas multi-tabla (dos de ellos diferidos) y cinco vistas de
  usuario. Los trece ajustes al modelo se documentaron de forma transparente.

  \item Se generaron cuatro escenarios de volumen reproducibles, hasta
  aproximadamente 3.86 millones de filas en el escenario de 1M, mediante un
  generador de datos sintéticos propio con semilla fija y carga por \texttt{COPY}.

  \item \textit{[Pendiente: agregar las conclusiones del experimento una vez
  obtenidos los tiempos reales. Deben mencionar qué consultas se beneficiaron
  del índice y en qué medida, el comportamiento de la consulta por rango (Q5),
  y el costo de los índices en escrituras.]}
\end{itemize}

% ==================================================================
\section{Anexos}

\subsection{Estructura del repositorio}

\begin{center}
\begin{tabular}{ll}
\toprule
\textbf{Archivo o carpeta} & \textbf{Contenido} \\
\midrule
\texttt{sql/01\_tables.sql}          & Tablas con todas sus restricciones. \\
\texttt{sql/02\_triggers.sql}        & Funciones y triggers. \\
\texttt{sql/03\_views.sql}           & Vistas de usuario. \\
\texttt{sql/04\_indexes.sql}         & Índices del experimento. \\
\texttt{sql/05\_drop\_indexes.sql}   & Quita los índices secundarios. \\
\texttt{sql/06\_experiments.sql}     & Consultas con \texttt{EXPLAIN (ANALYZE, BUFFERS)}. \\
\texttt{sql/07\_index\_overhead.sql} & Demostración del costo de los índices. \\
\texttt{querys/}                     & Las cinco consultas limpias (Q1--Q5). \\
\texttt{faker/}                      & Generador de datos sintéticos. \\
\texttt{dumps/}                      & Dumps restaurables de las cuatro bases. \\
\texttt{results/}                    & CSV de resultados y gráficos (pendiente). \\
\texttt{docs/}                       & Este informe y el detalle de cambios. \\
\bottomrule
\end{tabular}
\end{center}

\subsection{Video de la experimentación}

\textit{[Pendiente. El video debe mostrar la experimentación completa sobre la
base de 1 millón de registros: medición sin índices, creación de índices,
medición con índices y comparación de planes y tiempos. El guion detallado
está en \texttt{video/README\_guion.md} del repositorio.]}

\subsection{Planes de ejecución (Q5, base 1M)}

\textit{[Pendiente. Capturar la salida de \texttt{EXPLAIN (ANALYZE, BUFFERS)}
para la consulta Q5 sobre \texttt{bd\_literaria\_1m}, primero sin el índice
\texttt{idx\_material\_anio} y luego con él. Pegar los resultados aquí.]}

\end{document}
